\section*{Chapitre \ref{ch:b2:structure-determination} --- Détermination de structure~: RMN du \texorpdfstring{$^{13}$C}{13C}, SM et IR réunies}

\begin{solution}{exo:b2:structure-determination:1}
1,2-diméthylbenzène~: 4 (\ce{CH3}~; C1/C2~; C3/C6~; C4/C5). 1,3-~: 5 (\ce{CH3}~; C1/C3~; C2~; C4/C6~; C5).
1,4-~: 3 (\ce{CH3}~; C1/C4~; C2/C3/C5/C6).
\end{solution}

\begin{solution}{exo:b2:structure-determination:2}
\ce{CH3CH(OH)CH2CH3}. DEPT-135~: C1 (\ce{CH3}) vers le haut, C2 (CH) vers le haut, C3 (\ce{CH2}) vers le
bas, C4 (\ce{CH3}) vers le haut. DEPT-90~: C2 seul.
\end{solution}

\begin{solution}{exo:b2:structure-determination:3}
209~: carbonyle de cétone~; 170~: carbonyle d'ester (ou d'acide)~; 129~: carbone \omterm{def:b2:huckel:aromatic}{aromatique}~; 64~: carbone
lié à un oxygène~; 14~: \ce{CH3} alkyle.
\end{solution}

\begin{solution}{exo:b2:structure-determination:4}
$(2 \times 8 + 2 - 8)/2 = 5$~: un cycle benzénique (4) et un \ce{C=O} (1). Par exemple le benzoate de
méthyle, \ce{C6H5COOCH3}, ou l'acide 4-méthylbenzoïque.
\end{solution}

\begin{solution}{exo:b2:structure-determination:5}
Une insaturation, une cétone (\qty{1715}{cm^{-1}}, raie à 209, quaternaire)~; un \ce{CH2} et deux
\ce{CH3}, quatre carbones différents~: la butan-2-one, \ce{CH3COCH2CH3}.
\end{solution}

\begin{solution}{exo:b2:structure-determination:6}
Carbone~: la pentan-2-one présente cinq raies, la pentan-3-one, symétrique, trois. Protons~: la
pentan-2-one présente un singulet de trois protons (\ce{CH3CO})~; la pentan-3-one seulement un quadruplet
et un triplet. \omterm{def:b2:mass-spec-atomic:spectrum}{Spectre de masse}~: la pentan-2-one a son \omterm{def:b2:mass-spec-atomic:spectrum}{pic de base} à 43 (\ce{CH3CO+}) et un ion de
McLafferty à 58~; la pentan-3-one a son \omterm{def:b2:mass-spec-atomic:spectrum}{pic de base} à 57
(\ce{C2H5CO+}).
\end{solution}

\begin{solution}{exo:b2:structure-determination:7}
\ce{C4H8O2}, 88~: la paire quadruplet--triplet est un \ce{OCH2CH3} (le \ce{CH2} à 4.12 est sur l'oxygène),
le singulet un \ce{CH3CO}. C'est l'éthanoate d'éthyle, \ce{CH3COOCH2CH3}.
\end{solution}

\begin{solution}{exo:b2:structure-determination:8}
1,2-dichlorobenzène~: trois raies (C1/C2, C3/C6, C4/C5). 1,4-dichlorobenzène~: deux (C1/C4, les quatre
CH). La DEPT-90 montre deux raies de CH pour le premier, une pour le second.
\end{solution}

\begin{solution}{exo:b2:structure-determination:9}
C2 est un stéréocentre. Remplacer l'un ou l'autre des protons de C3 par un deutérium donne deux
diastéréoisomères, et non des énantiomères~: les deux protons sont diastéréotopes, dans des
environnements différents. Ils peuvent avoir des déplacements différents et se coupler entre eux comme
avec leurs voisins, si bien que le \ce{CH2} donne un multiplet complexe au lieu d'un simple quintuplet.
\end{solution}

\begin{solution}{exo:b2:structure-determination:10}
Cinq insaturations~; un cycle monosubstitué (trois raies de CH et une raie de C à 137.0, cinq H
\omterm{def:b2:huckel:aromatic}{aromatiques})~; une cétone (200.8, C)~; un groupe éthyle (\ce{CH2} à 31.8, quadruplet à 2.98 voisin du
carbonyle, \ce{CH3}
à 8.3). C'est la 1-phénylpropan-1-one, \ce{C6H5COCH2CH3}. La conjugaison avec le cycle délocalise les
électrons $\pi$ de \ce{C=O} et affaiblit la liaison~: la bande descend au-dessous des \qty{1715}{cm^{-1}}
d'une cétone saturée.
% ledger: ir:CO-ketone
\end{solution}

\begin{solution}{exo:b2:structure-determination:11}
Acide pentanoïque~: la très large bande \ce{O-H} de 2500 à \qty{3300}{cm^{-1}} et un proton vers
11--12~ppm. Butanoate de méthyle~: un singulet de trois protons dans le domaine des protons portés par
un carbone lié à un oxygène (3.3--4.5). Propanoate d'éthyle~: deux quadruplets, l'un sur l'oxygène,
l'autre voisin du carbonyle. Éthanoate de propyle~: un singulet de trois protons voisin du carbonyle
(2.0--2.4) et un triplet sur l'oxygène.
% ledger: ir:OH-acid, nmrr:acid, nmrr:alcohol, nmrr:ketone
\end{solution}

\begin{solution}{exo:b2:structure-determination:12}
Les deux raies disparaissent en \omterm{def:b2:structure-determination:dept}{DEPT}~: les deux carbones sont donc quaternaires, et la \omterm{def:b2:structure-determination:dept}{DEPT} seule ne peut
pas les distinguer. Les déplacements le peuvent~: les carbonyles d'ester se situent vers 167--171 et les
carbones \omterm{def:b2:huckel:aromatic}{aromatiques} entre 128 et 137 dans le diagramme de ce chapitre. Un carbone du cycle à 166.5 et un
carbonyle à 130.6 tomberaient tous deux loin de leurs classes~; l'attribution correcte est l'inverse.
\end{solution}

\begin{solution}{pb:b2:structure-determination:1}
\textbf{1.} $3.1/31.7 \approx 9.8~\%$~; $9.8/1.11 \approx 9$ carbones.
\textbf{2.} Masse paire~: pas d'azote (ou deux).
\textbf{3.} $150 - 108 = 42$~: \ce{H10O2} convient (10 + 32)~: \ce{C9H10O2}.
\textbf{4.} Dix carbones donneraient $M+1 \approx 11~\%$ de $M$, et non 9.8~\%~; de plus, l'IR montre un
carbonyle, qui exige un second oxygène en plus de la bande \ce{C-O}.
\textbf{5.} $(2 \times 9 + 2 - 10)/2 = 5$.
\textbf{6.} $108 + 10 \times 1.007825 + 2 \times 15.994915 \approx 150.0681$.
\textbf{7.} Une élongation \ce{C=O} à la position d'un ester saturé (vers \qty{1735}{cm^{-1}}).
\textbf{8.} L'élongation \ce{C-O} de l'ester.
\textbf{9.} Des liaisons \ce{C-H} \omterm{def:b2:huckel:aromatic}{aromatiques} (ou alcéniques).
\textbf{10.} Un groupe \ce{O-H}~: ni alcool ni acide carboxylique.
\textbf{11.} Non~: un carbonyle conjugué avec un cycle absorbe à plus faible nombre d'onde~; ici il se situe
là où absorbent les esters saturés, si bien qu'un groupe non conjugant le sépare du cycle.
\textbf{12.} 170.70~: \ce{C=O} d'ester~; 136.14~: carbone du cycle qui porte le substituant~; 128.56 et
128.24~: CH \omterm{def:b2:huckel:aromatic}{aromatiques}~; 66.24~: \ce{CH2} lié à l'oxygène~; 20.82~: \ce{CH3}.
\textbf{13.} 66.24, le \ce{CH2}.
\textbf{14.} 170.70 et 136.14, les deux carbones sans hydrogène.
\textbf{15.} 128.56 et 128.24.
\textbf{16.} Quatre~: le carbone qui porte le substituant, les deux CH ortho, les deux CH méta et le CH para.
On attend trois sortes de CH~; deux d'entre elles ont des déplacements trop proches pour être résolus à ce
champ, et partagent une raie.
\textbf{17.} Pas nécessairement~: hauteur n'est pas nombre (\cref{prop:b2:structure-determination:intensities}),
même si cette raie peut ici contenir deux sortes de CH superposées.
\textbf{18.} 7.33~: les cinq protons \omterm{def:b2:huckel:aromatic}{aromatiques}~; 5.09~: le \ce{OCH2}~; 2.06~: le \ce{CH3CO}.
\textbf{19.} Leurs atomes voisins (le carbone du cycle et l'oxygène pour le \ce{CH2}, le carbone du
carbonyle pour le \ce{CH3}) ne portent pas d'hydrogène.
\textbf{20.} Il est lié à l'oxygène électronégatif de l'ester et au cycle~: tous deux le déblindent.
\textbf{21.} \ce{C6H5CH2OCOCH3}, l'éthanoate de benzyle.
\textbf{22.} Dans l'isomère, le \ce{CH2} n'est pas lié à un oxygène, si bien qu'il ne pourrait pas apparaître à
5.09, et le méthyle, porté par l'oxygène, apparaîtrait dans le domaine 3.3--4.5, et non à 2.06 près d'un
carbonyle.
% ledger: nmrr:alcohol, nmrr:ketone
\textbf{23.} 108~: perte de 42 (cétène, \ce{CH2=C=O}) par réarrangement, le radical-cation de
\ce{C7H8O}~; 91~: \ce{C7H7+}, le cation benzyle~; 43~: \ce{CH3CO+}.
\textbf{24.} Sept sortes de carbone (\ce{C=O}, le carbone substitué du cycle, CH ortho, méta et para,
\ce{CH2}, \ce{CH3})~: \textbf{sept raies, dont une (celle du \ce{CH2}) pointe vers le bas} en DEPT-135.
\end{solution}
